\chapter{Trabajos futuros}
\label{cap:futuro}

Las líneas siguientes se derivan de las limitaciones del capítulo~\ref{cap:discusion}.

\begin{enumerate}
  \item \textbf{Niveles de comisión, latencia y volatilidad externa.} El simulador debe incorporar el intercambio con comisión de la ecuación~\eqref{eq:comision}, un retraso en la ejecución del arbitraje y una referencia con volatilidad, para convertir en resultados los factores que hoy son candidatos de diseño (tabla~\ref{tab:factores}). Con comisión, la banda de no arbitraje deja de ser un parámetro libre y queda determinada por $\gamma=-\log(1-f)$.
  \item \textbf{CEX con profundidad y latencia.} Sustituir la referencia exógena por un libro de órdenes que responda a las operaciones del AMM permitiría estudiar la realimentación bidireccional.
  \item \textbf{Liquidez concentrada.} Extender el AMM a posiciones por intervalos de precio, que cambian la relación entre reserva e impacto.
  \item \textbf{Reglas de proveedores y escala de órdenes.} Contrastar con datos de eventos reglas graduales o estimadas de respuesta de los proveedores y tamaños de orden ligados a las reservas; la sensibilidad actual fija reglas estilizadas y no identifica la conducta real de los proveedores.
  \item \textbf{Planificación de precisión.} Usar un piloto que replique el diseño de choques con signo y estratificar por dirección; el piloto original subestimó la varianza de H1.
  \item \textbf{Mapa reducido.} Probar mapas con más variables o con términos no lineales, y evaluar si mejoran al autorregresivo escalar; explorar otros criterios de clasificación de regímenes.
  \item \textbf{Validación de ingeniería con un par real.} Definir de antemano un par y un periodo; alinear la mejor compra/venta y profundidad del CEX con las reservas, intercambios y eventos de liquidez del DEX; registrar comisiones, gas, demora de ejecución, orientación de tokens, orden de bloques y derechos de uso de datos; después probar persistencia de recuperación, tamaño relativo de órdenes y desempeño reservado del modelo sustituto. No transferir directamente umbrales sintéticos a producción.
  \item \textbf{Evaluación de ingeniería y pedagogía.} Medir si el tablero mejora la capacidad de operadores o estudiantes para identificar recuperación transitoria y estrés relativo de órdenes/reservas; mantener esta evaluación separada de la validación de mercado.
\end{enumerate}
